\subsection{Finding Reach and Manipulation Capability After Failure}\label{sec:methods-reachability}
When an LMJ occurs, the bounds of the robot's reachable workspace change in intractable, non-linear ways. We must understand the new reachable area and discover what manipulation capabilities the robot retains across this failure-constrained region to continue to complete tasks.

To do this, we developed a reachability and ability analysis system, or \textit{Reach-Ability} (\cref{algo:reachability}). The task space of the robot is divided into small areas approximately the size of the Franka Emika Panda end-effector tip, $2$ cm (\cref{fig:methods}, Left Column, Top). For each discretized area, we check to see if the robot can find an inverse kinematics solution for prehensile actions. To be considered a valid prehensile pose, the end-effector must be in a range of orientations that allow for grasping, nominally between normal to the workspace surface and parallel to the workspace surface. To counter the effects of the randomness of the inverse kinematics solver, if no solution is found initially, we uniformly sample a perturbation from an $\epsilon$-ball inscribed within the discretization cell and attempt solving again. 

If there is still no solution after half the replan attempts are calculated, \texttt{changeInteraction()} (\cref{algo:reachability}, L9) will release the orientation constraints of PM manipulation and switch to NPM abilities.%change the interaction type from PM to NPM
Under the NPM abilities, the inverse kinematics solver (\cref{algo:reachability}, L10) will: 1) cycle through pre-selected interaction points on the robot: hand, wrist, and forearm and 2) relax its orientation error bound so that the solution is only restricted in position. If no solutions are found, then that point is considered unreachable. The map is smoothed through interpolation once the whole task space is explored (\cref{fig:methods}, Left Column, Bottom).  
We leverage this mapping of the robot's statically reachable, failure-constrained workspace to build a representation of the robot's action manifold by utilizing the notion of edge bundles that parameterize the set of all feasible actions at the intersection of the post-failure-accessible manifold and the motion primitive manifold. 
% This still doesn’t allow for the execution of actions which leads us to a hierarchical representation - a notion of edge bundles that parameterize the set of all feasible actions at the intersection of the failure manifold and the motion primitive manifold. Brute force iteration.

% an understanding of the robot's failure-constrained workspace, we can now construct the kinodynamic map within those bounds.

